\documentclass[10pt,a4paper]{amsart}
\usepackage[margin=19mm]{geometry}
\usepackage{amsmath,amssymb,mathtools}
\usepackage[scr=rsfs,cal=euler]{mathalfa}
\usepackage{xfrac}
\usepackage[unicode,hidelinks,bookmarks=false]{hyperref}
\newcommand{\ie}{i.e.\ }
\newcommand{\cf}{cf.\ }
\newcommand{\resp}{respectively\ }
\newcommand{\Subsection}{\subsection}
\providecommand{\nobreakdash}{\nobreak\hbox{-}}
\setlength{\emergencystretch}{2em}
\multlinegap0pt
\usepackage{tikz}
\usepackage{bm}
\usepackage{dsfont}
\usepackage{braket}
\usepackage{amssymb}
\usepackage{mathtools}
\usepackage{xcolor}
\newtheorem{cor}{Corollary}
\newtheorem{lem}{Lemma}
\newtheorem{prop}{Proposition}
\usepackage{cleveref}
\crefname{cor}{Corollary}{Corollarys}
\crefname{lem}{Lemma}{Lemmas}
\crefname{prop}{Proposition}{Propositions}
\title{Linear Combination of Hamiltonian Simulation with Commutator Scaling}
\author{Junaid Aftab, Dong An, Konstantina Trivisa}
\date{Source: arXiv:2606.11475v1 (CC BY 4.0)}
\begin{document}
\maketitle

\noindent\emph{Source and licence.} Linear Combination of Hamiltonian Simulation with Commutator Scaling. Version arXiv:2606.11475v1 (CC BY 4.0), distributed by arXiv under the Creative Commons Attribution 4.0 International licence, http://creativecommons.org/licenses/by/4.0/, as recorded in the arXiv licence metadata for this exact version and shown on the abstract page https://arxiv.org/abs/2606.11475v1. Full licence notice: \texttt{licenses/arxiv-2606.11475-CC-BY-4.0.txt}.\par\smallskip

\noindent\textit{Editorial scope.} Counted unit: the article's cor trap-budget (source0.tex lines 1407--1433). The article's complete original proof of it is delivered with it (source0.tex lines 1435--1481, one \texttt{proof} environment of 47 lines). No mathematics is written, repaired or re-derived: the statement and the proof are the article's own lines, in the article's own notation, and no retained line is substituted or re-flowed.  Retained uncounted as the source-local context the proof needs: lines 783--789; lines 1314--1316; lines 1329--1335; lines 1361--1371; lines 1384--1386.  Nothing was omitted from any retained span, so the package carries no omission record.  No retained line contains a citation, so the wrapper carries no bibliography and no bibliography entry.  No retained line contains a figure environment, an image-inclusion command or drawing code, so no image is delivered and the package declares no asset dependency.  Editorial formatting only, in addition: an AMS \texttt{amsart} preamble that restates the article's own declarations and macros used by the retained lines, namely the environments \texttt{cor}, \texttt{lem}, \texttt{prop} on separate counters, together with the standard packages those lines need.  The retained environments print in this document as Corollary 1, Lemma 1, Proposition 1 in order of appearance, whereas the article prints the same items with the numbering of its own full document.  The complete native source of the article as distributed in the arXiv e-print for this exact version is bundled unchanged as \texttt{sources/202552/source0.tex}.
\begin{cor}\label{radii-lower-bd}
Let \(\epsilon'>0\). The set \(\mathcal H_{\epsilon'}\) is non-empty for every \(\epsilon'>0\). For a kernel profile \(\vec{\theta}=(a,b,c,d)\), define
\begin{equation}
\vartheta(\vec\theta) := \frac{\sqrt{\pi}}{c}\left(\frac{b}{b+1}\right)^{a-1}\frac{\epsilon'}{e^{d-\frac{1}{4c^2}}}.
\end{equation}
If \(\vartheta(\vec\theta)<1\), then every \(R>2c\,\operatorname{erfc}^{-1}(\vartheta(\vec\theta))\) is in \(\mathcal H_{\epsilon'}\). If \(\vartheta(\vec\theta)\ge 1\), then every \(R>0\) is in \(\mathcal H_{\epsilon'}\).
\end{cor}

\begin{equation}\label{eq:rho-star-def}
\rho_\ast := \begin{cases} 1, & a=1,\\ \min\{1,b\}, & a>1. \end{cases}
\end{equation}

\begin{lem}\label{infinite-line-trapezoidal-estimate}
Let \(\rho>0\). Suppose that \(g\) is analytic on an open set containing \(\overline S_\rho\) and that \(g\) decays uniformly to zero on \(\overline S_\rho\) as \(|z|\to\infty\). If \(M>0\) satisfies $\int_{\mathbb R}\|g(x+i\beta)\|\,dx\le M$
for every \(\beta\in(-\rho,\rho)\), then, for every \(h>0\), we have
\begin{equation}\label{eq:trap-infinite-line}
\left\|\int_{\mathbb R}g(x)\,dx-h\sum_{q\in\mathbb Z}g(qh)\right\|\le \frac{2M}{e^{2\pi \rho/h}-1}.
\end{equation}
\end{lem}

\begin{prop}\label{uniform-lattice-quadrature}
Fix a kernel profile \(\vec{\theta}=(a,b,c,d)\), and let \(\rho\in(0,\rho_\ast)\), where \(\rho_\ast\) is defined in \cref{eq:rho-star-def}. Let \(R>0\), \(h>0\). Then
\begin{equation}\label{eq:uniform-lattice-error}
E_{\operatorname{quad}}(Q_{R,h}^{\operatorname{trap}}) = \|W_\infty-W_{R,h}^{\operatorname{trap},\operatorname{ideal}}\| \le E_{\operatorname{trap}}(h,\rho)+E_{\operatorname{tail},h}(R_h),
\end{equation}
where \(E_{\operatorname{trap}}(h,\rho):=\frac{2M_{a,b}(\rho)}{e^{2\pi \rho/h}-1}\). Moreover, we have $E_{\operatorname{tail},h}(R_h)\le E_{\operatorname{trunc}}(R_h)\le E_{\operatorname{trunc}}(R)$. 
Consequently, we have
\begin{equation}\label{eq:uniform-lattice-error-simple}
E_{\operatorname{quad}}(Q_{R,h}^{\operatorname{trap}})\le E_{\operatorname{trap}}(h,\rho)+E_{\operatorname{trunc}}(R).
\end{equation}
\end{prop}

\begin{equation}\label{eq:real-kernel-modulus}
|\hat f_{a,b}(k;c,d)| = \frac{(b+1)^{a-1}e^{d-\frac{1}{4c^2}}}{\sqrt{2\pi}}\frac{e^{-k^2/(4c^2)}}{\sqrt{1+k^2}(b^2+k^2)^{(a-1)/2}}.
\end{equation}

\begin{cor}\label{trap-budget}
Let \(\rho\in(0,\rho_\ast)\) and \(\epsilon_{\operatorname{mesh}},\epsilon_{\operatorname{tail}}>0\). Define
\begin{equation}\label{eq:tail-log-def}
\ell_{\operatorname{tail}} := \max\left\{1,\log\left(\frac{C_{\operatorname{tail}}}{\epsilon_{\operatorname{tail}}}\right)\right\}, \quad \ell_{\operatorname{mesh}} := \log\left(1+\frac{2M_{a,b}(\rho)}{\epsilon_{\operatorname{mesh}}}\right).
\end{equation}
Choose \(h=\min\left\{1,2c,\frac{2\pi\rho}{\ell_{\operatorname{mesh}}}\right\}\) and \(R=2c\sqrt{\ell_{\operatorname{tail}}}\), and set \(n_h=\lceil R/h\rceil\), \(R_h=n_hh\). Then $E_{\operatorname{quad}}(Q_{R,h}^{\operatorname{trap}})\le \epsilon_{\operatorname{mesh}}+\epsilon_{\operatorname{tail}}$. Moreover, we have
\begin{align}
|\mathcal I_{Q_{R,h}^{\operatorname{trap}}}| &= 2n_h+1 = \mathcal O\left(1+(1+c)\sqrt{\ell_{\operatorname{tail}}}+\frac{c}{\rho}\sqrt{\ell_{\operatorname{tail}}}\,\ell_{\operatorname{mesh}}\right), \label{eq:trap-comparison-cardinality-1}\\
R_{Q_{R,h}^{\operatorname{trap}}} &= R_h = \mathcal O\left(c\sqrt{\ell_{\operatorname{tail}}}+\min\left\{1,2c,\frac{\rho}{\ell_{\operatorname{mesh}}}\right\}\right). \label{eq:trap-comparison-cardinality-2}
\end{align}
For a fixed admissible kernel profile and balanced budgets \(\epsilon_{\operatorname{mesh}}\asymp\epsilon_{\operatorname{tail}}\asymp\epsilon_q\), we obtain
\begin{equation}\label{eq:trap-balanced-cardinality}
|\mathcal I_{Q_{R,h}^{\operatorname{trap}}}| = \mathcal O\left(\frac{c}{\rho}\log^{3/2}\left(\frac1{\epsilon_q}\right)\right), \quad R_{Q_{R,h}^{\operatorname{trap}}} = \mathcal O\left(c\sqrt{\log\left(\frac1{\epsilon_q}\right)}\right),
\end{equation}
for fixed \(\rho\) and \(c\). For \(\nu\ge0\), define
\begin{equation}\label{eq:trap-sharp-moment-scale}
\mathfrak M_{\nu,a,b}(c) := \begin{cases} 1, & \nu<a-1,\\ 1+\log(1+c), & \nu=a-1,\\ 1+c^{\nu-a+1}, & \nu>a-1. \end{cases}
\end{equation}
Then there exists a constant \(C_{\nu,a,b}^{\operatorname{trap}}>0\), depending only on \(\nu,a,b\), such that
\begin{equation}\label{eq:trap-sharp-moment-bound}
M_{\nu,R,h}^{\operatorname{trap}}\le \widetilde C_{\operatorname{ker}}C_{\nu,a,b}^{\operatorname{trap}}\mathfrak M_{\nu,a,b}(c).
\end{equation}
In particular, \(\alpha_{R,h}^{\operatorname{trap}}=M_{0,R,h}^{\operatorname{trap}}\le \widetilde C_{\operatorname{ker}}C_{0,a,b}^{\operatorname{trap}}\mathfrak M_{0,a,b}(c)\). Equivalently, the trapezoidal moments satisfy
\begin{equation}\label{eq:trap-sharp-moment-asymptotic}
M_{\nu,R,h}^{\operatorname{trap}} = \begin{cases} \mathcal O(\widetilde C_{\operatorname{ker}}), & \nu<a-1,\\ \mathcal O(\widetilde C_{\operatorname{ker}}\log(1+c)), & \nu=a-1,\\ \mathcal O(\widetilde C_{\operatorname{ker}}c^{\nu-a+1}), & \nu>a-1. \end{cases}
\end{equation}
\end{cor}

\begin{proof}
The choice of \(h\) controls the mesh error. By \cref{infinite-line-trapezoidal-estimate}, \(E_{\operatorname{trap}}(h,\rho)\le \epsilon_{\operatorname{mesh}}\). It remains to control the truncation error. By \cref{radii-lower-bd}, \(E_{\operatorname{trunc}}(R)\le C_{\operatorname{tail}}\operatorname{erfc}(R/(2c))\). Since \(\operatorname{erfc}(x)\le e^{-x^2}\), the condition \(R\ge 2c\sqrt{\ell_{\operatorname{tail}}}\) implies \(E_{\operatorname{trunc}}(R)\le \epsilon_{\operatorname{tail}}\). Combining these estimates with \cref{uniform-lattice-quadrature} gives
\begin{equation}
E_{\operatorname{quad}}(Q_{R,h}^{\operatorname{trap}})\le E_{\operatorname{trap}}(h,\rho)+E_{\operatorname{trunc}}(R)\le \epsilon_{\operatorname{mesh}}+\epsilon_{\operatorname{tail}}.
\end{equation}
Since \(n_h=\lceil R/h\rceil\), we have \(|\mathcal I_{Q_{R,h}^{\operatorname{trap}}}|=2n_h+1\le 2R/h+3\). Moreover,
\begin{equation}
\frac1h = \max\left\{1,\frac1{2c},\frac{\ell_{\operatorname{mesh}}}{2\pi\rho}\right\} = \mathcal O\left(1+\frac1c+\frac{\ell_{\operatorname{mesh}}}{\rho}\right).
\end{equation}
Substituting \(R=2c\sqrt{\ell_{\operatorname{tail}}}\) gives the bound on \(|\mathcal I_{Q_{R,h}^{\operatorname{trap}}}|\) in \cref{eq:trap-comparison-cardinality-1}. Moreover, \(R_h=n_hh\le R+h\), giving the bound on \(R_{Q_{R,h}^{\operatorname{trap}}}\) in \cref{eq:trap-comparison-cardinality-2}. The balanced-budget estimates follow from \(\ell_{\operatorname{mesh}}=\ell_{\operatorname{tail}}=\mathcal O(\log(1/\epsilon_q))\), for fixed \(\rho\) and \(c\). It remains to bound the LCU normalization and discrete moments. For real \(k\), \cref{eq:real-kernel-modulus} implies
\begin{equation}\label{eq:kernel-sharp-majorant}
\frac{1}{\sqrt{2\pi}}|\hat f_{a,b}(k;c,d)| = \widetilde C_{\operatorname{ker}}\frac{e^{-k^2/(4c^2)}}{\sqrt{1+k^2}(b^2+k^2)^{(a-1)/2}}.
\end{equation}
Therefore,
\begin{equation}\label{eq:trap-moment-sharp-start}
M_{\nu,R,h}^{\operatorname{trap}}\le \widetilde C_{\operatorname{ker}}h\sum_{q\in\mathbb Z}\frac{e^{-(qh)^2/(4c^2)}|qh|^\nu}{\sqrt{1+(qh)^2}(b^2+(qh)^2)^{(a-1)/2}}.
\end{equation}
For \(k>0\), we have
\begin{equation}\label{eq:trap-denom-min-bound}
\frac{1}{\sqrt{1+k^2}(b^2+k^2)^{(a-1)/2}}\le \min\left\{\frac{1}{b^{a-1}},\frac{1}{k^a}\right\}.
\end{equation}
Indeed, the first bound follows from \(\sqrt{1+k^2}\ge1\) and \((b^2+k^2)^{(a-1)/2}\ge b^{a-1}\), while the second follows from \(\sqrt{1+k^2}\ge k\) and \((b^2+k^2)^{(a-1)/2}\ge k^{a-1}\). Let
\begin{equation}\label{eq:trap-kappa-b-def}
\kappa_b := \begin{cases} 1, & a=1,\\ b^{(a-1)/a}, & a>1. \end{cases}
\end{equation}
We split the full-line sum in \cref{eq:trap-moment-sharp-start} into the regions \(|qh|\le \kappa_b\) and \(|qh|>\kappa_b\). Since \(h\le1\), the first contribution is bounded by a constant depending only on \(\nu,a,b\). More precisely,
\begin{equation}
h\sum_{|qh|\le\kappa_b}\frac{e^{-(qh)^2/(4c^2)}|qh|^\nu}{\sqrt{1+(qh)^2}(b^2+(qh)^2)^{(a-1)/2}}\le C_{\nu,a,b}.
\end{equation}
On the complementary region, \cref{eq:trap-denom-min-bound} gives
\begin{equation}
h\sum_{|qh|>\kappa_b}\frac{e^{-(qh)^2/(4c^2)}|qh|^\nu}{\sqrt{1+(qh)^2}(b^2+(qh)^2)^{(a-1)/2}}\le h\sum_{|qh|>\kappa_b}e^{-(qh)^2/(4c^2)}|qh|^{\nu-a}.
\end{equation}
The latter Riemann sum has three regimes. If \(\nu<a-1\), then \(|k|^{\nu-a}\) is integrable at infinity, and hence
\begin{equation}
h\sum_{|qh|>\kappa_b}e^{-(qh)^2/(4c^2)}|qh|^{\nu-a}\le C_{\nu,a,b}.
\end{equation}
If \(\nu=a-1\), then the summand behaves like \(1/|qh|\) until the Gaussian cutoff scale \(|qh|\sim c\), and therefore
\begin{equation}
h\sum_{|qh|>\kappa_b}e^{-(qh)^2/(4c^2)}|qh|^{-1}\le C_{a,b}\bigl(1+\log(1+c)\bigr).
\end{equation}
If \(\nu>a-1\), comparison with the corresponding Gaussian integral gives
\begin{align}
h\sum_{|qh|>\kappa_b}e^{-(qh)^2/(4c^2)}|qh|^{\nu-a} &\le C_{\nu,a,b}+C_\nu\int_0^\infty e^{-k^2/(4c^2)}k^{\nu-a}\,dk \le C_{\nu,a,b}+C_\nu c^{\nu-a+1}.
\end{align}
Combining these three cases with \cref{eq:trap-moment-sharp-start} proves \cref{eq:trap-sharp-moment-bound}. The normalization bound is the case \(\nu=0\), and the asymptotic statement follows directly from the definition of \(\mathfrak M_{\nu,a,b}(c)\).
\end{proof}

\end{document}
